%! Unit Opener: Why Algebra Fluency Matters — Unit 1, Lesson 1.0 — Warm-Up
\documentclass[10pt]{article}
\usepackage{atda-article}
\usepackage{atda-boxes}
\newcommand{\TallMath}[1]{$\displaystyle #1\rule[-1.4em]{0pt}{3.2em}$}
\setlength{\workrowsep}{6pt}   % handwriting room; moves blank and key together

\begin{document}

\pageheader{Unit 1, Lesson 1.0}{Warm-Up --- Unit 1 Diagnostic}
% NAMESTRIP: no \namedateperiod here. The student writes their name once, on the
% cover this component is stapled behind. See references/conventions.md.

\vspace{0.15in}

\begin{notesbox}{Where Are You Starting? --- Five Skills Unit 1 Builds On}
\small
This is a \textbf{diagnostic, not a quiz}. Each item checks one skill this unit rebuilds. Show
your work; we score it together and you will record the results on your unit roadmap.

\begin{enumerate}[leftmargin=1.6em, itemsep=4pt, topsep=4pt]

\item \textbf{Exponent rules.} Simplify: $(3x^2y)(-4x^5y^3)$
\begin{work}
  (3x^2y)(-4x^5y^3) &= (3)(-4)\,x^{2+5}\,y^{1+3} \\
                    &= -12x^7y^4
\end{work}

\item \textbf{Multiplying binomials.} Multiply and write in standard form: $(x+3)(x-7)$
\begin{work}
  (x+3)(x-7) &= x^2-7x+3x-21 \\
             &= x^2-4x-21
\end{work}

\item \textbf{Greatest common factor.} Factor out the GCF: $12a^3-18a^2$
\begin{work}
  12a^3-18a^2 &= 6a^2(2a)-6a^2(3) \\
              &= 6a^2(2a-3)
\end{work}

\item \textbf{Solving a linear equation.} Solve for $x$: $5(x-2)=3x+4$
\begin{work}
  5x-10 &= 3x+4 \\
     2x &= 14 \\
      x &= 7
\end{work}

\item \textbf{Evaluating in context.} A fireworks shell is launched from a $64$-foot cliff. Its
height above the water, in feet, $t$ seconds after launch is $h(t)=-16t^2+48t+64$.
Find $h(2)$, then write one sentence saying what that number means for the shell.
\begin{work}
  h(2) &= -16(2)^2+48(2)+64 \\
       &= -64+96+64 \\
       &= 96
\end{work}
\writeline

\end{enumerate}
\end{notesbox}

\end{document}
